\documentclass[10pt,letterpaper]{article}
\usepackage{cornell-notes}

\title{Introduction}
\author{}
\date{}
\setDocTitle{Introduction}
\setDocSubtitle{ISO/IEC/IEEE 42010:2022}
\setDocOwner{ISO 42010 Cornell Notes}
\setDocDate{\today}
\setCornellCollection{ISO/IEC/IEEE 42010:2022 Cornell Notes}
\setCornellUnitType{Introduction}
\setCornellUnitNumber{Introduction}
\setCornellUnitTitle{Introduction}
\hypersetup{pdftitle={Introduction},pdfsubject={Cornell notes},pdfauthor={}}

\begin{document}
\maketitle
\begin{CornellOverviewBox}{Overview}
{\Large\bfseries\textcolor{CNavy}{Introduction: Introduction}}\\[3pt]
{\small\textbf{Classification:} Introductory}\\
{\small\textbf{Learning objective:} Explain the role of introduction and connect it to related architecture-description concepts.}
\end{CornellOverviewBox}
\vspace{3pt}

\begin{CornellNotesTable}
\CornellNoteRow{Purpose / key concept}{The complexity of human-made entities has grown to an unprecedented level.}
\CornellNoteRow{Core premise}{Architecture descriptions are tangible work products used to communicate and cooperate around the otherwise intangible, abstract architecture of an entity.}
\CornellNoteRow{Supporting resources}{Architecture description frameworks codify conventions and common practices; architecture description languages codify ways of expressing architectures within communities and domains.}
\CornellNoteRow{Normative language}{SHALL denotes a requirement; SHOULD denotes a recommendation; MAY denotes a permission.}
\CornellNoteRow{Scope of use}{Architecture descriptions support design, development, analysis, evaluation, maintenance, risk mitigation, planning, decision support, review, training, trade studies, and lifecycle activities.}
\CornellNoteRow{Example / application}{of entities include the following: Enterprise, organization, solution, system (including software systems), subsystem, process, business, data (as a data item or data structure), application, information technology (as a collection), mission, product, service, software item, hardware item, product line, family of systems, system of systems, collection of systems, collection of applications.}
\CornellNoteRow{Interpretive note}{\begin{itemize}[leftmargin=*,nosep,topsep=1pt]
\item ISO/IEC/IEEE 42020 specifies a set of processes for architecting which can be employed in support of creating one or more ADs. The architecture elaboration process in ISO/IEC/IEEE 42020 is especially relevant for creation of ADs.
\end{itemize}}
\CornellNoteRow{Active recall}{\begin{itemize}[leftmargin=*,nosep,topsep=1pt]
\item What is the central role of Introduction?
\item How does Introduction connect to adjacent architecture-description concepts?
\end{itemize}}
\end{CornellNotesTable}


\begin{CornellSummaryBox}{Summary}
The complexity of human-made entities has grown to an unprecedented level.
\end{CornellSummaryBox}

\end{document}
